\chapter{第十一讲\quad 犯罪形态}
\label{lecture:11}
\noindent{\small\hyperref[ref:lecture:11]{查阅本讲扩展规则}}\par

对同一行为人、同一候选罪名和同一犯罪过程 \(e\)，令 \(S(e)\) 为已经着手，\(K(e)\) 为该罪既遂要件全部满足，\(X(e)\) 为由于意志以外原因未能进展到既遂，\(V(e)\) 为在犯罪过程中自动停止。这里的时间单位是能独立评价的犯罪过程。计划仍在实施时，不能因为一次击打落空就立即给整个过程贴上终局标签。

\section{第一节\quad 犯罪预备}

\paragraph{法条和模型}
《刑法》第22条把为犯罪准备工具或制造条件纳入预备。令 \(I\) 为准备行为指向特定故意犯罪，\(P\) 为实际准备工具或制造条件，\(X_0\) 为外部原因阻止着手。教材用于判定终局预备的规则是
\[
I\land P\land\neg S\land X_0
\Longrightarrow \text{犯罪预备的终局形态}.
\]
仅有犯意表露而没有准备行为，不满足 \(P\)。为犯罪而实施的动作自身还可能符合另一罪名的既遂，须独立检验。例如非法取得枪支可能涉及独立枪支犯罪，不能用预备标签覆盖该事实。预备犯的从宽是可以比照既遂从轻、减轻或者免除。

\paragraph{推导}
甲为在夜间抢劫乙准备工具并到约定地等待，还未对乙施加足以压制反抗的暴力或胁迫，即被发现并控制。准备工具和选定场所给出 \(I\land P\)，未逼近实行阶段给出 \(\neg S\)，被控制给出 \(X_0\)，合取后可认定抢劫预备。只在聊天时说想抢钱而未付诸准备，推不出 \(P\)。保持上述条件但把被控制改成有条件继续、甲自愿停止，不能推出 \(X_0\)，应转入预备阶段的中止审查。

\section{第二节\quad 犯罪未遂}

\paragraph{法条和模型}
《刑法》第23条。就具备故意及其他必要条件的犯罪过程，
\[
S\land\neg K\land X\Longrightarrow \text{犯罪未遂}.
\]
教材用行为是否已向具体罪状的实行推进、是否使法益面临现实危险来审查 \(S\)。抢劫一般检查是否开始压制反抗，诈骗检查是否已经向相对人实行相应欺骗，保险诈骗注意开始索赔与仅制造事故的区别。隔离犯、间接正犯的着手另有学理争议，应说明采用的判断基准。未遂犯可以比照既遂从轻或者减轻处罚。

\paragraph{推导}
甲为夺取乙的财物，持棍逼乙交钱。乙抵抗，甲尚未取得财物，亦未达到教材所述抢劫既遂的相应伤害条件，随即被制止。甲开始以威胁压制反抗，给出 \(S\)。相关既遂要件未满足，给出 \(\neg K\)。外部制止使甲无法继续，给出 \(X\)。三项合取，得抢劫未遂。若甲仅藏在场地附近等候，\(S\) 不成立，不能沿该式认定未遂。

\paragraph{不能犯的边界}
用 \(D_m(b)\) 表示在危险判断模型 \(m\) 下，行为 \(b\) 具有法益侵害危险。教材对照主观说、事后客观危险说与行为时具体危险说。主观说重行为人判断，事后说重查明的客观情况，具体危险说从行为时一般人的处境观察法益危险。对事后发现剂量不足的案件，先固定对象、实际手段和当时可知信息，再比较 \(D_m\)，不能用 \(\neg K\) 直接推出 \(\neg D_m\)。例如甲把空地上的稻草人认作乙，场地当时确无他人处在射击范围内。在排除现实危险的教材模型中，不能仅由杀人意图推出可罚的杀人未遂。若甲误认乙而实际射向另一在场者，则对象错认并没有消除对人的危险，须另论故意及归责。

\section{第三节\quad 犯罪中止}

\paragraph{法条和模型}
《刑法》第24条。令 \(C\) 为停止时犯罪过程仍未终结，\(K\) 为同一犯罪已达既遂，\(E\) 为先前行为已经足以使目标结果自行进展，\(Q\) 为行为人有效防止应归属于前行为的目标结果，\(D\) 为先前犯罪过程造成依法应计入中止处罚的损害。按对有效性要求实际阻止作用的模型，
\[
\begin{aligned}
C\land\neg K\land V\land\neg E\land\text{放弃继续实行}
    &\Longrightarrow \text{中止},\\
C\land\neg K\land V\land E\land Q
    &\Longrightarrow \text{中止},\\
\text{中止}\land\neg D &\Longrightarrow \text{应当免除处罚},\\
\text{中止}\land D &\Longrightarrow \text{应当减轻处罚}.
\end{aligned}
\]
\(E\) 依具体行为与结果进展判断，不能仅凭作案工具是否用完。继续犯可能已既遂而尚未终了，故 \(C\) 与 \(\neg K\) 须分别判断。自动性的核心是行为人认为仍能继续却选择停止。若误以为已经被堵截、无法继续而停止，不能仅凭客观上尚有退路推出 \(V\)。

\paragraph{推导}
甲故意杀乙并造成足以致死的伤势，在结果发生前主动及时送乙就医，救治阻止了可归属于刺击的死亡。致命伤给出 \(E\)，有继续放任结果的可能而主动实施救治给出 \(C\land V\)，救治有效给出 \(Q\)。故在既遂结果未发生时成立故意杀人中止。伤势已构成重伤并可归于先前刺击，\(D\) 为真，适用应当减轻。若甲离去、邻人独立救人，虽同样未死亡，甲没有实施中止措施，不可从 \(\neg K\) 单独推出 \(Q\) 或 \(V\)。

\paragraph{有效性分歧与边界}
教材区分措施的可能有效性与结果归属上的实际有效性。措施明显无用或仅敷衍，不因恰好被第三人救活就自动变为中止。行为人已采取合理且真诚的救助措施，另有第三人更早救助时，教材在特定事实下认可中止。此时应明示所用解释，分别陈述行为人已做的救助与真正阻止结果的因果链。若施救途中出现独立的新风险并导致死亡，先判断原犯罪行为的风险是否实现。若死亡应归于先前犯罪，不能仅凭救助意愿认定中止。若死亡由独立风险所致，另审原犯罪形态与新风险的责任。目标结果和中止处罚中的损害 \(D\) 不等同。

\section{第四节\quad 犯罪既遂}

\paragraph{教材模型}
既遂判据依罪名设定。令 \(A_g\) 为罪 \(g\) 的实行行为，\(R_g\) 为该罪要求的结果或法定危险状态，\(C_g\) 为法律所要求的归责关联，\(J_g\) 为其余要件。对需要结果归责的罪，
\[
K_g\ \Longleftrightarrow\ A_g\land R_g\land C_g\land J_g.
\]
有的罪以完成特定行为或达到法定危险作为既遂，不可把所有 \(R_g\) 都定义成实害。对同一人、同一罪和同一终局过程，既遂时 \(K_g=1\)，未遂、中止和终局预备均要求该既遂未实现。因此同一过程不能既遂又未遂。数个独立过程可以分别定形态，须先说明过程为何分离。

\paragraph{推导}
甲假称在售货，乙识破，按侦查安排给甲一笔钱。假话与交付各自存在。以 \(E_f\) 表示受骗后的错误认识，以 \(D_f\) 表示因错误而处分，则乙识破给出 \(\neg E_f\)，也就不能建立欺骗经由 \(E_f,D_f\) 产生财产损害的因果链。若甲已经着手、具有非法占有目的，且相应的未遂可罚条件成立，甲因意志以外原因未得逞，可认定诈骗未遂。不能从钱到了甲手中直接推出既遂。相同检验可以迁移到敲诈勒索中的恐惧处分和抢劫中的压制反抗。

\paragraph{终局反例}
故意杀人中一击未中，行为人仍能继续却当场放弃，整体过程可进入中止。相反，行为被彻底制止、第一次过程已结束后，数日再起犯意实施犯罪，须另检验第二次的阶段和结果。既遂后退回财物或补救可进入量刑评价，不倒推消除既遂。

\section*{怎么记与闭卷自测}
先圈人、罪和过程，再画一条时间轴。\textbf{预看备和拦，遂看手和外，止看能而停，成看罪状齐。} 逐题复述五步。先问罪状有没有齐，再问已经着手没有，再问为何停止，再问是否还需有效防止结果，最后把中止造成的其他损害单独列出。记忆时用同一案件依次替换三个条件，练习从未着手到已着手、由外力到自愿、由仅放弃到有效救助。
\begin{enumerate}
\item 甲已开始抢劫乙，第一次夺取落空，仍有机会继续而自动放弃，没有造成损害。应认何种形态？\par
答：同一过程尚可继续，落空并未终结。若符合其他条件，成立中止，应当免除处罚。
\item 甲投毒后打电话自首，但未通知任何能及时救助的人，乙因毒物死亡。是否构成中止？\par
答：可归责死亡已发生，既遂成立。自首另按其规则评价，不能当作有效防止结果。
\item 行为人实施欺骗并收到转账，付款人并未受骗。还须查什么？\par
答：对既遂须查错误认识与处分、损害间的归责链。既遂不成立时，再查着手、外部原因和该罪未遂的可罚条件。
\end{enumerate}
